% prcsid $Id: guisdoc.tex 1.30 Thu, 30 Dec 2004 14:14:05 +0100 basile $
% prcsproj $ProjectHeader: Guis 1.62 Thu, 30 Dec 2004 14:14:05 +0100 basile $
\documentclass[a4paper,11pt]{article}
\usepackage{times}
\usepackage{hevea}
\usepackage{url}
\usepackage{hyperref}
\usepackage{lastpage}
%\usepackage{fancyhdr}
\usepackage{epsf}
\usepackage[latin1]{inputenc}
%\pagestyle{fancy}
\title{GUIS - a GUI widget server\\ {\small \it
% $Format: "release $ReleaseVersion$ on $ProjectDate$\\\\"$
release 1.6 on Thu, 30 Dec 2004 14:14:05 +0100\\
{\tiny \tt
% $Format: "prcsproj $ProjectVersion$"$
prcsproj 1.62
}
}}
\author{Basile STARYNKEVITCH -- \mailto{basile\_NOSPAM@starynkevitch.net.invalid}\\
\url{http://www.starynkevitch.net/Basile/index_en.html}\\
\small{8, rue de la Fa�encerie, 92340 Bourg La Reine, France}}

\begin{document}

\newenvironment{guisscript}{\begin{quote}%
%HEVEA\navy
\begin{latexonly}\small\end{latexonly}
\tt}{\end{quote}}

\newenvironment{applcode}{\begin{quote}%
%HEVEA\green
\begin{latexonly}\small\end{latexonly}
\tt}{\end{quote}}

\newenvironment{guisrequest}{\begin{quote}%
%HEVEA\purple
\begin{latexonly}\small\end{latexonly}
\tt}{\end{quote}}

\newenvironment{guisreply}{\begin{quote}%
%HEVEA\blue
\tt}{\end{quote}}

\newcommand{\figimagebox}[1]{%
%HEVEA{\imgsrc{#1.jpeg}}
\begin{latexonly}\epsfbox{#1.eps}\end{latexonly}%
}


\maketitle
\newpage
\tableofcontents
\newpage

\begin{center}
\begin{tabular}{|r|l|c|}
\hline
\multicolumn{3}{|c|}{\large Software Description}\\
\hline
\textbf{Name} & {GUIS}\\
\textbf{License} & {GNU General Public License}\\
\textbf{Author} & {Basile {\sc Starynkevitch}}\\
\textbf{Version} &
                                %$Format: "{$ReleaseVersion$}\\\\"$
{1.6}\\
\textbf{Development system} & {Linux/Debian/Sid x86}\\
\textbf{Programming Language} & {C}\\
\hline
\multicolumn{3}{|c|}{\large Software Dependencies}\\
\hline
GTK  2.4 (or 2.2) \\ and related libraries \\
{\small (Glib, Pango, Atk)} & required \\
Python 2.3 (or 2.2) & recommended \\
PyGTK 2.2 & recommended \\
Ruby 1.8 & recommended \\
Ryby-Gnome2 0.11.x & recommended \\
Slang 1.4.x & optional \\
Slgtk 0.5.x & optional \\
\hline
\end{tabular}
\end{center}


Please be nice to send me (\mailto{basile@starynkevitch.net}) an email
if you use this information and this Guis software.

Guis is available (as a gnuzipped source tarball) from
% $Format:"\\url{http://www.starynkevitch.net/Basile/guis-$ReleaseVersion$.tar.gz}"$
\url{http://www.starynkevitch.net/Basile/guis-1.6.tar.gz}
%
and this document is on
\url{http://www.starynkevitch.net/Basile/guisdoc.html}. See also my
home page on \url{http://www.starynkevitch.net/Basile/} or Guis page
\url{http://freshmeat.net/projects/guis/} on Freshmeat for
announcement of newer versions. Please feel free to send suggestions,
patches, criticisms, etc...

\section{Overview and usage}

This section gives a short overview with the classical adder example.
Then the usage details are given.


\subsection{Motivations and related stuff}

{\em Guis} is a small widget server. It is a gtk2 (see
\url{http://www.gtk.org/}) based program listening on a pipe for
widget requests (requests are Python [or Ruby] scripts - see
\url{http://www.python.org/} using the PyGTK 2.0 binding of GTK2 to Python
- see \url{http://www.daa.com.au/~james/software/pygtk/}) and
outputting events or replies {\em Guis} is useful for programs {\small
  (in particular, setuid programs or (ruby,ocaml,perl...) scripts)}
which do not want to link in a full widget toolkit but prefer to
delegate the user interface to another process.

The choice of the scripting language is not critical provided it does
have a full Gtk2 binding. Porting Guis to another scripting language
should be easy.

{\small (Many years ago, I was a satisfied user of Sun OpenWindows
  system with its programmable NeWS [widget] server. I still miss that
  widget server (see a message I posted in october 1993 on
  \url{http://www.zendo.com/vsta/mail/1/0146.html} which is copied on
  \url{http://starynkevitch.net/Basile/NeWS_descr_oct_1993.html}) I don't
  understead why the NeWS team, which also probably designed Java, did
  not consider to embed the JVM inside the X11 server, as a standard
  X11 extension, and embedding a toolkit inside Java, like Swing does,
  is not an answer.)}

A similar project is {\em entity} - see \url{http://www.entity.cx/}.
The major difference between {\em entity} and {\em guis} is that {\em
  guis} is a server (listening for orders on a pipe...) while {\em
  entity} is a script engine.

The IRAF widget server - see
\url{http://iraf.noao.edu/iraf/web/projects/x11iraf/} had similar goals.
And PicoGui \url{http://picogui.org/} is also server based.

The XUL system of mozilla (see \url{http://www.xulplanet.com/}) also
describes an interface with XML.

The (previously Berlin, now) Fresco server should be a corba server
for widgets - which seems nearly dormant.  see
\url{http://www.fresco.org/}



\subsection{Introduction}

{\em Guis} is a graphical user interface program communication with a
client application (using separate protocols). The application send
widget building requests to {\em Guis} (so these requests are input
for {\em Guis}) and handles widget events sent from {\em Guis}.
Usually Guis is started with a small Python initial script which
defines common functions and build some widgets. The requests are
Python source code chunks. The replies (i.e. events sent back from
{\em Guis} to the application) are just some textual lines sent (by
some Python code calling {\tt guis\_send}).

Actually, {\em Guis} is strongly dependent on GTK2, and depends less
of Python. The code is designed to make {\em Guis} easily
portable\footnote{To port Guis to another language you just have to
  link {\tt gguis.c} with a file similar to {\tt py\_guis.c} for your
  scripting language which provides the following functions : {\tt
    guis\_initialize\_interpreter(void)} (called once to initialise the
  interpreter), {\tt guis\_load\_initial\_script(char* scriptname)} (called
  once with either the initial script name -a file path- or NULL),
  {\tt guis\_interpret\_request(char* request)} (called for every request),
  and {\tt guis\_end\_of\_input\_hook(int timeout)} (called at end of input
  with a timeout in milliseconds). These last 3 functions should
  return 0 if successful, or a static C string describing the error.}
to any other scripting interpreter able to evaluate requests in
textual strings, provided this interpreter has a binding to GTK2. The
only reason I use Python here is the availability of a nearly complete
binding of Python to GTK2. (I would prefer some other scripting
languages). To remind that {\em Guis} is using Python its binary is
called {\tt pyguis}.

Since version 1.3, {\em Guis} is also interfaced to Ruby. See section
\ref{sec:rubyport} below.

\subsubsection{callbacks}

The initial script (or the application) is responsible for installing
appropriate callbacks with the {\bf connect} primitive (or equivalent)
of the scripting language (Python or Ruby). 

{\Large IMPORTANT}
Callbacks in Guis should be robust: {\bf callbacks cannot raise
  uncaught exceptions} (because they are run by the {\tt
  gtk\_main\_loop} in {\em Guis}, outside of the (Python or Ruby...)
interpreter. Applications should encapsulate callbacks with the
appropriate mechanism (catch ... ) in the scripts.


\subsubsection{initial script}

Usually, {\em Guis} runs an initial script (in Python or Ruby) which
is interpreted by the scripting language before entering the {\tt
  gtk\_main\_loop}. This initial script usually builds the widget and
defines some application specific functions (to implement the protocol
specific to your application).

The initial script is run once. It is specified with the {\tt -s}
option or with the {\tt -scripter} trick. See below \ref{invocation}
and the man page.

You might (if possible) end your initial script with a call to {\tt
  gtk\_main\_loop} as provided by the (Python or Ruby) GTK2 binding.
Calling it will make the request evaluation better under control and
might permit exceptions in callbacks (see your documentation of the
GTK binding you are using). Then you have to exit explicitly {\em
  Guis} (by calling the {\tt exit} primitive of Python or Ruby) or
  tell {\em Guis} (using {\tt\verb+guis_main_loop_in_script+})
  that you call the main loop.

I still strongly advise against uncaught exceptions in callbacks.



\subsubsection{protocols}

Every request sent from the application to {\em Guis} should end with
two consecutive newline\footnote{The newline is coded decimal 10 in
  ASCII or IsoLatin1} characters {coded in C as \tt \verb+\n\n+} or
with a formfeed (coded in C as {\tt \verb+\f+}, decimal 12).

Obviously requests cannot contain (inside) a double newline which is a
convention suitable for most scripting languages (including Python,
Ruby, Ocaml, Lua, Rep-Lisp, Slang, ...).

A convenient way to debug {\em Guis} initial scripts is to run {\tt
  pyguis} with an explicit FIFO input: make it with {\tt mkfifo
  /tmp/fifo} and then run {\tt pyguis -s yourscript -i /tmp/fifo -o
  -}; in another xterm, run {\tt cat >> /tmp/fifo} and don't forget to
end every request with a double newline (ie return) character.

Events or replies sent from {\em Guis} to the application are single
lines (maybe very long) ended with a newline. They should not contain
any control character (eg newline or formfeed) inside. Requests and
replies are asynchronous (a request can be sent without any replies
and vice versa).

The driving idea of {\em Guis} is that the input and output protocols are
tailored to your application. On the input side (requests from
your application to {\em Guis}) the protocol is usually made of calls
to specific functions defined in the initial script. On the output
side (replies or events from {\em Guis} to your application) the
protocol is defined by sending (thru an appropriate primitive, usually
{\tt guis\_send}, of the scripting language) arbitrary lines to your
application from callbacks.

\subsubsection{other toolkits (Qt3)?}

It would be interesting to have a similar approach with QT3. I tried,
and leave some (bad, incomplete, not even compilable) C++ code under
the {\tt bad\_qt\_stuff/} directory of this {\em Guis}. Feel free to
reuse this code (under a GNU license). My main problem was lack of
good binding to QT3 and threading problems (notably threads are nearly
incompatible with an embedded Python).
\subsection{Small example}

For illustration purpose, suppose we have an application which
computes the sum of 2 integers, and we want to give it a nice
graphical interface containing two (editable) textual entry widgets, a
quit button, and a label widget displaying the sum.

\begin{figure}[h]
\begin{center}
\figimagebox{imgdemo}

\caption{Simple example demo window}
\end{center}
\end{figure}

\subsubsection{protocols}

We have to think first about the messages sent from the application to
{\em Guis}. We need first to {\tt start} the interface (giving some
nice title). We will need to display a sum using {\tt displaysum} and
to display an error message using {\tt displayerror}. And
we need to stop the demo, thru a {\tt stopdemo} function. All these
functions are Python functions defined in the initial script file.

We also need to define the messages sent from {\em Guis} back to the
application. We will send a plain {\tt END} for end, and a more
complex message starting with {\tt ADD} to ask the application to make
an addition (displaying the result with a {\tt displaysum} request.
The {\tt ADD} message should contain the textual content of the two
entry widgets. Since the textual content can be anything (it could
even contain control characters like newlines) it should be encoded.
We use a C like encoding convention, so will usually send {\tt ADD "1"
  "3"} -or even {\tt \verb!ADD "\t1" "3"!} if the first entry starts
with a tab\footnote{How to enter a tab character inside a Gtk entry
  widget is left as an exercise to the reader.}. The application is in
charge of checking that the entries contain valid numbers.

The {\em Guis} server program buffers all input (python requests) and
output (event replies), reading and writing as soon as possible.

A typical exchange between the application and Guis might be as
follow; first the application starts and sends

\begin{guisrequest}
\verb+start("pid 1234")+
\end{guisrequest}

Then {\em Guis} shows the window and let the user interact with it. Some
user interaction makes {\em Guis} send back messages like

\begin{guisreply}
ADD "2" "5"
\end{guisreply}


To which the application responds with
\begin{guisrequest}
\verb!displaysum(2,5,7)!
\end{guisrequest}

When the user closes the window, {\em Guis} send back 
\begin{guisreply}
END
\end{guisreply}

To which the application responds with
\begin{guisrequest}
stopdemo()
\end{guisrequest}

and then exits.


When run with the {\tt -T} flag, {\em Guis} opens a window to show the
trace of all requests and replies (this is useful for script
and application debugging):

\begin{figure}[h]
\begin{center}
\figimagebox{imgtrace}

\caption{Protocol trace window}
\end{center}
\end{figure}

\subsubsection{initial Python script}

We write a small Python initial script. A special trick in {\em Guis}
is that if {\tt Guis} is invoked with a name (i.e. {\tt argv[0]} in C
parlance) ending with {\tt -scripter} then the next (second) argument
is the initial script name. Hence we can start our script with

\begin{guisscript}
\verb+#! /usr/bin/env pyguis-scripter+\\
\verb+# file guisdemo_script in -*- python -*-+
\end{guisscript}

With such a trick, our initial script can be invoked by any {\tt
  pyguis-scripter} found in our {\tt \$PATH}. We make it a symbolic
link to the {\tt pyguis} executable.

We need to tell python to use the {\tt gtk} module (provided by pygtk)
and the {\tt guis} module (builtin inside {\tt pyguis}).

\begin{guisscript}
\verb+import gtk+\\
\verb+import guis+
\end{guisscript}

Next, we need to define a callback used by the {\it quit} button; it
just sends back the {\tt END} string to the application

\begin{guisscript}\begin{verbatim}
def end_cb(*args):
    guis.guis_send("END")
\end{verbatim}\end{guisscript}

We also define a callback used when text entries are updated. It uses
the {\tt guis.to\_guis} primitive to convert a C string to its textual
representation but we could have used Python {\tt repr} function.

We need to define the {\tt start} function, invoked by the application
in its first request, to build the graphical widgets and connect them
to callbacks. It first builds a window and its contained vertical box
(using GTK2 calls in Python):

\begin{guisscript}\begin{verbatim}
def start(welcomsg):
    global window, xent, yent, sumlab
    window = gtk.Window(gtk.WINDOW_TOPLEVEL)
    window.set_title("Guis Demo")
    vbox = gtk.VBox(gtk.FALSE,2)
    window.add(vbox)
\end{verbatim}\end{guisscript}

Then it builds the other widgets (details deleted here, see the source
of {\tt guisdemo\_script} file). At last, it makes the {\it quit}
button, connect it to the {\tt end\_cb} callback, add it into {\tt
  vbox} and show all of the {\tt window}:

\begin{guisscript}\begin{verbatim}
    # ...
    quitbut = gtk.Button("quit")
    quitbut.connect("clicked", end_cb)
    vbox.add(quitbut)
    window.show_all()
\end{verbatim}\end{guisscript}

We need to define the {\tt displaysum} function

\begin{guisscript}\begin{verbatim}
def displaysum(x,y,sum):
    sumlab.set_markup(('<i>%d</i> + <i>%d</i>' % (x,y))
                       + (' = <big>%d</big>' % sum))
\end{verbatim}\end{guisscript}
    
We need to define the {\tt displayerror} function. To avoid
messing the Gtk2 (pango provided) XML-like markup, we convert the
message to its XML representation (i.e. using {\tt \verb+&lt;+}
for {\tt <} etc...) using the {\tt guis.xml\_coded} primitive.

\begin{guisscript}\begin{verbatim}
def displayerror(message):
    sumlab.set_markup('<b>ERROR:</b> '
                      + (guis.xml_coded(message)))
\end{verbatim}\end{guisscript}
    
A {\tt stopdemo} function is also needed (see the source file).

\subsubsection{client application}

We suppose the client application is written in C. You can code it in
any language able to communicate on channels in a textual way. We
comment here parts of the file {\tt guisdemo\_client.c}. You don't need
to link any {\em Guis} specific library to it!

We declare a big line buffer, and the requests and replies files. We
could also use the {\tt glibc} specific {\tt getline} function which
dynamically allocates the line buffer.

\begin{applcode}\begin{verbatim}
  char linbuf[1024];
  FILE *toguis = stdout;
  FILE *fromguis = stdin;
\end{verbatim}\end{applcode}
    
  
  
At first, we want to send a request like
{\tt \verb+start("<b>guis</b> demo")+}. Requests may start with a
comment used to help identify them in error messages\footnote{Actually
  requests are identified by their first line in error messages.}. 
\begin{applcode}\begin{verbatim}
  fprintf (toguis, "#initial start\n"
           "start(\"pid %d\")\n\n",
           (int) getpid ());
  fflush (toguis);
\end{verbatim}\end{applcode}

Never forget to flush your request channel very often, and to end
every request with two newlines.

Of course we need a loop to read events (or replies) messages from
{\em Guis} - each of them is a single (sometimes very long) line ended
with a single newline.

\begin{applcode}\begin{verbatim}
  while (!feof (fromguis)) {
    fgets (linbuf, sizeof (linbuf) - 1, fromguis);
\end{verbatim}\end{applcode}
    
if the reply line starts with {\tt ADD} we scan it appropriately and ask
to display a fancy line like ``{\it 2} + {\it 3} = {\large 5}''
otherwise (bad scan because of non-numeric entries) we display
``{\bf invalid input}''

\begin{applcode}\begin{verbatim}
    if (!strncmp (linbuf, "ADD", 3)) {
      int x=0,  y=0,  pos=0;
      if (sscanf (linbuf, "ADD \"%d\" \"%d\" %n", 
                  &x, &y, &pos) > 0 && pos > 0) {
        fprintf (toguis, "#good sum\n"
                 "displaysum(%d,%d,%d)\n\n",
                 x, y, x + y);
      } else {
        fprintf (toguis, "#bad input\n"
                 "displayerror(\"bad input\")\n\n");
      }
\end{verbatim}\end{applcode}

If the reply is {\tt END} we stop gently (by sending a {\tt stopdemo()}
request and exiting): 
\begin{applcode}\begin{verbatim}
    } else if (!strncmp (linbuf, "END", 3)) {
      fprintf (toguis, "#stop\n" "stopdemo()\n\n");
      fflush (toguis);
      sleep (1);
      exit (0);
    } 
\end{verbatim}\end{applcode}

After warning against unexpected input lines, we flush the request
channel and end the loop.
\begin{applcode}\begin{verbatim}
    fflush (toguis);
  }; // end of while feof
\end{verbatim}\end{applcode}

Normally the while loop should never be ended, since our guis python
script should signal termination with {\tt END} (handled above).


\section{Reference}

\subsection{installing Guis}

You need Python (2.2.x or 2.3.y from \url{http://www.python.org/}),
GTK (2.2 or better from \url{http://www.gtk.org/}) and PyGTK (1.99.18
or 2.0 or better from
\url{http://www.daa.com.au/~james/software/pygtk/}) to build {\tt
  pyguis} (the Python version of {\em Guis}). You need Ruby (1.8.x)
from \url{http://www.ruby-lang.org/} and ruby-gnome 0.9.1 or later from
\url{http://ruby-gnome2.sourceforge.jp/} to build {\tt ruguis} (the
Ruby version of {\em Guis}). I built a thread-less
Python-2.2{\footnote{I passed the -disable-threads option to Python's
    configure script} which works with Guis. I am using GNU gcc (3.3)
  and GNU make (3.80).  You may add a local {\tt \_local.mk} file
  containing definitions for your installation, such as {\tt
    PREFIX=/usr}, {\tt CC=gcc-3.3}, {\tt
    PYTHONCFLAGS=-I/usr/include/python2.2} or {\tt
    PYTHONLDFLAGS=-lpython2.2} {\tt RUBY=ruby} etc...  you may even
  edit the Makefile.

First configure, either with {\tt make config} or with the {\tt
  ./Configure} script. Run it with {\tt -help} to get usage
  information.

Then run {\tt make} then {\tt make install} (which usually requires to be
root). 

You may build only the Ruby version with {\tt make ruguis} or only the
Python version with {\tt make pyguis}.


To run the demo, you probably need to add {\tt .} or the {\em Guis}
source directory to your {\tt \$PATH} before running {\tt
  guis\_demo.sh} or you can run the {\tt guisdemo\_script -p
  guisdemo\_client -T} command. use {\tt guisdemo\_rubyscript} to run
the Ruby version.


\subsection{invoking Guis}
\label{invocation}

See the man page in the source distribution for complete reference of
invocation, or invoke the binay with the {\tt --help} option.

{\em Guis} is usually invoked as {\tt pyguis} command, or indirectly
as \tt pyguis-scripter} if started by its initial script.

Guis can be started in a slave fashion (after the application has
started) by specifing its input and output channels (thru the {\tt -i}
and {\tt -o} options). You may specify file descriptors or paths as
channels.

Guis can also be started as a master, by giving the application
command as argument with {\tt -p}. You usually need to quote this
argument (because of your shell) unless the command has no spaces!

The {\tt -T} option is interesting to show the exchange between {\em
  guis} and its application in a separate window. This is very useful
  to debug your initial script or your application. 

The {\tt -D} option (disabled with {\tt -DNDEBUG} compile flag) show
lots of debugging information (to debug {\tt pyguis} itself).

The {\tt -L} option writes all requests into a log file.

the {\tt -I} or {\tt --input-enconding} option set the input encoding
(as supported by GLib2 on input channels). Likewise for  {\tt -O} or {\tt --output-enconding} 

\section{Guis for Python}

\subsection{invocation}

as usual, see section \ref{invocation} above.

\subsubsection{added Python primitives}

In alphabetical order, here are the names wired in the {\tt guis}
builtin Python module. You need of course to learn and use the modules
provided by {\tt pygtk} to actually build any GTK2 widget! You should
not explicitly call {\tt gtk\_main\_loop} from Python, since it is is
already called by {\tt pyguis}

\subsubsection{end\_of\_input\_hook}

With a callable argument, sets the hook called at end of input. Always
return the previous hook (even without any arguments). You probably
also need to set the end timeout using {\tt end\_timeout} below.

\subsubsection{end\_timeout}

Get (without arguments) or set (with an integer argument) the timeout
in milliseconds after which {\tt pyguis} exits at end of input. Useful
with {\tt end\_of\_input\_hook} above.

\subsubsection{guis\_send}

Send the string argument on the output channel to the application. The
string should not contain control characters (this is not checked)
otherwise the application might have trouble scanning it. A newline is
added if needed. (You may use many Python packages to build the sent
string; eg you could send XML stuff).

\subsubsection{main\_loop\_in\_script}

Get (without argument) or set (with a truth-value argument) the flag
telling {\em Guis} that the main loop {\tt gtk\_main} was called from
the initial script.

\subsubsection{nb\_replies}

Get (without arguments) the number of sent replies (event lines sent
to application). You might set this number with an integer argument
(but I see no reason to do this).

\subsubsection{nb\_requests}

Get the number of processed Python requests. You might set this number
with an integer argument (but I see no reason to do this).

\subsubsection{pipe\_check\_period}

Get (without arguments) or set (with an integer argument) the period
in milliseconds (should be 0 or 50 to 10000) to check for application
termination (in master mode).  Useful with {\tt end\_of\_input\_hook}
above.

\subsubsection{to\_guis}
\label{pyth:toguis}
Convert an object (or many of them, considered as a tuple) to a string
sendable to the application using the following algorithm:

\begin{itemize}
  
\item if the object is string, represent it in C syntax (except that
  double quotes are escaped as {\tt \verb+\Q+} so they are only string
  delimiters in the converted string) , so the string made of 4
  characters (a, tabulation, double-quote, z) is converted to the 7
  character string {\tt \verb+"a\t\Qz"+}

\item if the object is an integer, represent it in decimal notation
  
\item if the object is a double, represent it like in C with a {\tt
    \#} prefix\footnote{Distinguishing floating numbers from integers
    with a prefix should make parsing easier for the application.}, eg
  {\tt \#3.14}
  
\item if the object is a tuple, convert each component using the same
  {\tt to\_guis} builtin function and catenate each substring with
  spaces for separation - if any component fails to be converted, the
  entire tuple conversion fails.

\item if the object has a {\tt to\_guis} attribute, fetch it: if it is
  a string, return it, if it is callable, apply it to the object.

\item if the object has a {\tt to\_guis} method, call it (without any
  arguments except the recieving object).

\item otherwise, fail

\end{itemize}


\subsubsection{xml\_coded}

Quickly convert a Python string or unicode to its XML representation
(so the 3 characters string {\tt \verb+a<b+} is converted to the 6
characters string {\tt \verb+a&lt;b+}), escaping characters per XML
requirements:

\begin{itemize}
\item {\tt \verb+&+} as {\tt \verb+&amp;+} 
\item {\tt \verb+'+} as {\tt \verb+&apos;+}
\item {\tt \verb+"+} as {\tt \verb+&quot;+} 
\item {\tt \verb+<+} as {\tt \verb+&lt;+} 
\item {\tt \verb+>+} as {\tt \verb+&gt;+} 
\item other strict ASCII printable characters are kept identical
\item all other characters including control and IsoLatin1 accentuated
  characters like {\tt � �} are represented by a numerical entity
  (character number in decimal) like 
{\tt \verb+&#231;+} for the lowercase c with cedilla {\tt \verb+�+}
\end{itemize}

The result of {\tt guis.xml\_coded} is a string with only ASCII
printable characters (coded 32-126).


\section{Guis with Ruby}
\label{sec:rubyport}

{\em Guis} has been ported to Ruby 1.8.
See \url{http://www.ruby-lang.org/} for more on Ruby. The binary name
is {\tt ruguis} (including the {\tt ruguis-scripter} trick) and has
the same invocation as the {\tt pyguis} (Python) version. The Ruby
port is in the file {\tt ru\_gguis.c}

This port uses the {\it ruby-gnome2} binding of GTK2. See
\url{http://ruby-gnome2.sourceforge.jp/}.

The demo works in a Ruby way - it is in the {\tt guisdemo\_rubyscript}
file and can be run as
%
{\tt \verb+ruguis -T -D -s guisdemo_rubyscript -p guisdemo_client+} 
%
(provided that {\tt .} is inside your \verb+$PATH+)




\subsection{open questions}



I don't know if a builtin module can have virtual variables (as
defined by the {\tt rb\_define\_virtual\_variable} C function of the
Ruby runtime).

I would like to print more (e.g. the current environment in {\tt
  debug\_extra}) but I don't know how to code it.


\subsection{Ruby API}

There is a builtin {\tt Guis} Ruby module. It contains the {\tt
  guis\_send} primitive to send a string back to the application.

{\tt \$guis\_nbreq}, {\tt \$guis\_nbsend} and {\tt \$guis\_pipecheckperiod}
are (global) virtual variables (implemented thru {\tt
  rb\_define\_virtual\_variable} C function).

{\tt \$guis\_main\_loop\_in\_script} is a global variable which when
set to a true value avoid calling {\tt gtk\_main} after interpreting
the initial script (which should hence call {\tt Gtk.main}
explicitly).

Conversion to XML notation is done by the {\tt to\_xml} method added
to the existing {\tt String} class.

Conversion to a C-like notation (like {\tt to\_guis} in Python above
\ref{pyth:toguis}) is done by the {\tt to\_guis} method with built-in
implementations for arrays, strings, integers, floats, symbols. The
user could add other methods to existing classes by Ruby code like


\begin{guisscript}\begin{verbatim}
ExistingClass.class_eval { 
  def to_guis
    return "result"
  end 
}
\end{verbatim}\end{guisscript}


Then end of input hook is settable with

\begin{guisscript}\begin{verbatim}
Guis::on_end_of_input do |timeout|
  # user end input hook
done
\end{verbatim}\end{guisscript}

To remove the end of input hook, just do
\begin{guisscript}\begin{verbatim}
#removing end of input hook
Guis::on_end_of_input
\end{verbatim}\end{guisscript}


\section{experimental and incomplete Slang version of GUIS using slgtk}

See \url{http://s-lang.org/} for the Slang interpreter. See
\url{http://space.mit.edu/cxc/software/slang/modules/slgtk} for the
Slgtk binding of Slang to GTK.

The GUIS port to Slang is incomplete and barely tested. See the source
code for details. Feedback is welcome.

\section{Feedback}

\subsection{Feedback welcome}

Please send comments, criticisms, suggestions and patches to
\mailto{basile\_NOSPAM@starynkevitch.net.invalid} (but remove the {\tt
  \_NOSPAM} and {\tt .invalid} from the email) mentioning {\tt guis}
in the subject line. Feel free to suggestion new features. Tell me
about any success stories (ie Guis use) in your applications!

\begin{rawhtml}
    
  <p>I am interested by your <a href=
  "../showinterest.php?name=guisdoc">feedback</a> <small>(to this
  <tt>guisdoc</tt> page)</small></p>

 <form action="../interest.php" method="post">
  Did you find this page:<br>

  <input type="hidden" name="page" value="guisdoc"> 

  <input type="hidden" name="rev" value="$Revision: 1.30 $"> 

  <input type="radio" name="rate" value="1">Highly Interesting
      <small>(1)</small>; &nbsp; 

  <input type="radio" name="rate" value="2">Interesting 
       <small>(2)</small>; &nbsp; 

  <input type="radio" name="rate" value="3">Uninteresting
      <small>(3)</small>; &nbsp; 

  <input type="radio" name="rate" value="4">Awefully Uninteresting 
       <small>(4)</small><br>
       
   <input type="submit" value="Send feedback.">
 </form>

\end{rawhtml}


 \subsection{Change log}

\begin{description}

\item[version 1.6] (december 30, 2004) - minor code clean.

\item[version 1.5] (may 10, 2004) - ported to Gtk2.2, added Slang
  preliminary port (and non-working port to Perl).

\item[version 1.4] (september 5, 2003) bugfix (read buffer uses {\tt
    GString}) and added {\tt\verb+guis_main_loop_in_script+}),
  with both Python2.2 and Python2.3 support.

\item[version 1.3] (september 2, 2003) contains the experimental Ruby port
(and warns against exception in callbacks). 

{\small 
\begin{itemize}
\item 1.3.1 added word wrap in
trace window, 
\item 1.3.2 added logfiles with {\tt -L} with the Gtk.init
requirement in user scripts for Ruby
\item 1.3.3 {\scriptsize (september 3, 2002)} shows the versions info
  in trace window
\item 1.3.4 {\scriptsize (september 4, 2002)} removed the Gtk.init
  requirement in the user initial script
\end{itemize}}

\item[version 1.2] (august 31, 2003) minor fixes: just added the {\tt
  guis.xml\_coded} primitive and updated the demo and the
documentation.

\item[version 1.1] (august 30, 2003) made a major switch to Python using
PyGTK previous versions used Lua).

\item[versions $<= 0.3$]  (march 22, 2003) and older was Lua based
  {\small (with a Gtk binding to Lua generated by a CommonLisp script)}

\end{description}

\end{document}

